% ---
% id: receiver/rem:kamin
% kind: rem
% title: "Relation to a published multiclick weight bound"
% summary: >-
%   A published construction bounds a photon-number subspace weight from an
%   observed multiclick rate, for one row per detector. [[thm:multiclick-bound]] is
%   instead an operator inequality for every input state and ports of any rank;
%   no step here depends on the citation.
% prominence: standard
% status: accepted
% uses: [thm:multiclick-bound]
% assumes: []
% refs:
%   - "kamin2024improved, Thm. 1, Eq. (58)"
% provenance: newly-derived
% lean:
% ---

\begin{remark}\label{rem:kamin}
A published route to a related quantity is
\cite[Thm.~1, Eq.~(58)]{kamin2024improved}, which writes a multiclick POVM
element as $M^{\mathrm{mult}}(G)=\id-\dyad{0}-\sum_iF_{\mathrm{sic}}$ for a
passive linear-optical setup with mode transformation $G$, and then bounds the
weight of a photon-number subspace using an \emph{observed} multiclick rate
together with constants obtained by partitioning the rows of $G$. That is a
different object from \eqref{eq:mcbound}: it is a statement about a measured
statistic and a subspace weight, whereas Theorem~\ref{thm:multiclick-bound} is an
operator inequality on each photon-number block that holds for every input
state and is used as a device fact rather than as an observation. The
construction is also less general in the direction that matters here: it is
written for a one-row-per-detector transformation, whereas
Theorem~\ref{thm:multiclick-bound} is proved for arbitrary port Grams, hence for ports
of arbitrary rank. The citation is for orientation, and no step of this
contract depends on it.
\end{remark}
